\scanpage{356}
\mspage{95}

Cela provient de l'inclusion
\[
\underline{\mathrm{Hom}}(\tau,\tau')\subset\underline{\mathrm{Hom}}^{\varepsilon}(I_\tau,I_{\tau'})
\]
et de la description des catégories $\underline{\mathrm{Ch}}_\tau,\underline{\mathrm{Ch}}_{\tau'}$ comme catégories de foncteurs, telle
\[
\underline{\mathrm{Ch}}_{\Sigma,\tau}(X)\simeq\underline{\mathrm{Hom}}_{\Sigma_\tau,\Sigma}(I_\tau,X),
\]
où $\Sigma_\tau$ désigne l'ens. des flèches dans $I_\tau$ qui sont soit \deleted{\unsure{des iso}} inversibles, soit \uline{rétrogrades}. (Cet ens. de flèches est admissible, car stable par composition.) On a désigné par $\underline{\mathrm{Hom}}_{\Sigma_\tau,\Sigma}$ la sous-catégorie pleine de $\underline{\mathrm{Hom}}$ formée des foncteurs qui appliquent $\Sigma_\tau$ dans $\Sigma$. Or de façon générale, si
\[
(I,\Sigma_I)(I',\Sigma_{I'}),\quad(X,\Sigma)
\]
sont trois catégories, munies d'ens. admissibles de flèches, \deleted{\unsure{on trouve les foncteurs suivants}} \inserted{considérant les foncteurs} $\varphi:I\mathrel{\unsure{-|\to}}I'$ \deleted{\unsure{un foncteur}} avec $\varphi(\Sigma_I)\subset\Sigma_{I'}$, alors le foncteur \deleted{\unsure{f}} \unsure{usuel|usuelle|naturel}
\[
\underline{\mathrm{Hom}}(I,I')\longrightarrow\underline{\mathrm{Hom}}\bigl(\underline{\mathrm{Hom}}(I',X),\underline{\mathrm{Hom}}(I,X)\bigr)
\]
(transposé du foncteur de composition
\[
\begin{aligned}
\underline{\mathrm{Hom}}(I,I')\times\underline{\mathrm{Hom}}(I',X)&\longrightarrow\mathrm{Hom}(I,X))\\
(\varphi,F')&\longmapsto F'\circ\varphi),
\end{aligned}
\]
\deleted{\unsure{induit}} \inserted{se généralise en} un foncteur
\[
\underline{\mathrm{Hom}}_{\Sigma_I,\Sigma_{I'}}(I,I')\longrightarrow
\underline{\mathrm{Hom}}\bigl(\underline{\mathrm{Hom}}_{\Sigma_{I'},\Sigma}(I',X),\underline{\mathrm{Hom}}_{\Sigma_I,\Sigma}(I,X)\bigr),
\tag*{($\unsure{*|**}$)}
\]
transposé du foncteur de composition
\[
\underline{\mathrm{Hom}}_{\Sigma_I,\Sigma_{I'}}(I,I')\times\underline{\mathrm{Hom}}_{\Sigma_I,\Sigma}(I',X)\mathrel{\unsure{-|\to}}\underline{\mathrm{Hom}}_{\Sigma_I,\Sigma}(I,X)
\]
induit par $(*)$. \deleted{\unsure{Comme tout}}

% archive sheet number: 355
